\documentclass[journal]{IEEEtran}
\IEEEoverridecommandlockouts

\usepackage[T1]{fontenc}
\usepackage{newtxtext}
\usepackage{newtxmath}
\usepackage{microtype}
\usepackage{cite}
\usepackage{amsmath,amssymb}
\usepackage{graphicx}
\usepackage{booktabs}
\usepackage{multirow}
\usepackage{array}
\usepackage{xcolor}
\usepackage{url}
\usepackage{balance}
\usepackage{dblfloatfix}
\usepackage{hyperref}
\hypersetup{colorlinks=true,linkcolor=black,citecolor=black,urlcolor=black,allcolors=black}

% Float placement parameters
\renewcommand{\topfraction}{0.95}
\renewcommand{\bottomfraction}{0.95}
\renewcommand{\textfraction}{0.05}
\renewcommand{\floatpagefraction}{0.80}
\renewcommand{\dbltopfraction}{0.95}
\renewcommand{\dblfloatpagefraction}{0.80}

\setlength{\textfloatsep}{6pt plus 1pt minus 1pt}
\setlength{\floatsep}{6pt plus 1pt minus 1pt}
\setlength{\dbltextfloatsep}{6pt plus 1pt minus 1pt}
\setlength{\dblfloatsep}{6pt plus 1pt minus 1pt}

% Supplementary section and table numbering
\renewcommand{\thetable}{S.\Roman{table}}
\renewcommand{\thefigure}{S.\arabic{figure}}
\renewcommand{\thesection}{S.\Roman{section}}
\renewcommand{\thesubsection}{S.\Roman{section}.\Alph{subsection}}

\begin{document}

\title{Supplementary Material: SafeAdapt: Geometry-Preserving Optimal Transport and Fail-Closed Safeguards for Test-Time Adaptation in Digital Pathology}

% Submission Mode Toggle: \blindtrue for double-blind review, \blindfalse for camera-ready
\newif\ifblind
\blindfalse

\ifblind
\author{Anonymous~Authors%
\thanks{Supplementary Material submitted for double-blind peer review.}}
\else
\author{Avinash Reddy Cheerapareddy, Gunnam Sanjana\\
Department of Computer Science and Engineering,\\
SRM University-AP, Amaravati, Andhra Pradesh, India.\\
avinashreddy\_cheerapareddy@srmap.edu.in, sanjana\_gunnam@srmap.edu.in%
\thanks{Avinash Reddy Cheerapareddy and Gunnam Sanjana are with the Department of Computer Science and Engineering, SRM University-AP, Amaravati, Andhra Pradesh, India (e-mail: avinashreddy\_cheerapareddy@srmap.edu.in; sanjana\_gunnam@srmap.edu.in). Supplementary Material for IEEE Transactions on Medical Imaging (TMI). Source code and benchmark configurations: \url{https://github.com/Avirech7676/UTTA-Med}.}}
\fi

\maketitle

\section{Complete Per-Seed Target Evaluation Metrics}
\label{sec:supp_per_seed}

Table~\ref{tab:supp_perseed_full} presents the complete evaluation results across all five random seeds ($\{42, 123, 2024, 7, 99\}$) on Target Hospital~2 ($N=85{,}054$ patches across 10 whole-slide images) for all evaluated test-time adaptation methods and control baselines. SafeAdapt-OT consistently outperforms unconstrained entropy minimization (Tent) and existing gated baselines across discrimination (AUROC), classification accuracy ($F_1$ score), and calibration (ECE).

\section{Empirical Stress-Testing: Streaming Batch Size, Benign Slides, and Weakly-Supervised WSI-MIL}
\label{sec:supp_empirical_defenses}

\subsection{Batch-Size Streaming Robustness ($B=1$ to $B=64$ with FIFO Memory Queue)}
In streaming edge inference or whole-slide scanner feeds, tissue patches may arrive in variable-sized micro-batches. Standard Sinkhorn-Knopp Optimal Transport requires sample ensembles to construct empirical row marginals. To resolve this boundary, SafeAdapt-OT incorporates a fixed-capacity FIFO feature memory queue ($Q=64$ patches) when streaming mini-batch size $B < 16$. 

As demonstrated in Table~\ref{tab:supp_batch_size_sweep}, single-patch streaming ($B=1$) with FIFO memory buffering achieves an AUROC of $0.9482 \pm 0.0062$ and F1 score of $0.8398$, within $0.0013$ of standard batching ($B=32$, $0.9495$ AUROC), with maximum parameter displacement $\|\Delta\theta\|_\infty \le 0.0384$ strictly contained within the certified $L_\infty \le 0.05$ ball. This provides empirical proof that Sinkhorn prototype transport is not constrained to large batch sizes.

% Advance counter for Table S.II
\addtocounter{table}{1}
\begin{table}[b]
\centering
\caption{Batch Size Sensitivity \& Streaming FIFO Queue Performance ($n=5$ Seeds).}
\label{tab:supp_batch_size_sweep}
\resizebox{\columnwidth}{!}{%
\begin{tabular}{@{}ccccc@{}}
\toprule
Batch Size ($B$) & FIFO Queue ($Q=64$) & AUROC & F1 Score & Max Drift $\|\Delta\theta\|_\infty$ \\
\midrule
1 & Yes & 0.9482 & 0.8398 & 0.0384 \\
2 & Yes & 0.9485 & 0.8405 & 0.0410 \\
4 & Yes & 0.9489 & 0.8411 & 0.0425 \\
8 & Yes & 0.9491 & 0.8418 & 0.0441 \\
16 & No & 0.9493 & 0.8420 & 0.0450 \\
32 & No & \textbf{0.9495} & \textbf{0.8423} & 0.0458 \\
64 & No & 0.9496 & 0.8425 & 0.0460 \\
\bottomrule
\end{tabular}%
}
\end{table}

\subsection{Resilience on Purely Benign Slides (0.0\% Tumor Burden)}
A critical clinical challenge in unsupervised adaptation is handling completely benign slides ($0.0\%$ tumor presence). In unconstrained entropy minimization (Tent), attempting to minimize Shannon entropy on benign tissue forces the network into an attractor state, producing a catastrophic false-positive rate of $42.6\%$ and erroneously escalating benign lymph nodes to AJCC pN1a macro-metastases. 

In contrast, under SafeAdapt-OT, Theorem~1's Lipschitz bound strictly clamps parameter drift ($\Delta p \le 0.0460$), while selective entropy gating skips uninformative confident stroma. As shown in Table~\ref{tab:supp_benign_slides}, SafeAdapt-OT maintains a low false positive rate of $1.1\%$, completely avoiding attractor collapse and correctly preserving pN0 nodal staging.

\begin{table}[b]
\centering
\caption{Model Behavior on 100\% Benign Slides (0.0\% True Tumor Burden).}
\label{tab:supp_benign_slides}
\resizebox{\columnwidth}{!}{%
\begin{tabular}{@{}lcccc@{}}
\toprule
Method & True Tumor & False Pos Rate & Predicted Stage & Collapse Status \\
\midrule
Frozen Source & 0.0\% & 0.8\% & pN0 (Correct) & No \\
Tent (ICLR 2021) & 0.0\% & 42.6\% & pN1a (Fatal Error) & YES (Attractor) \\
EATA (ICML 2022) & 0.0\% & 14.2\% & pN0(i+) (Discordant) & Partial \\
SAR (ICLR 2023) & 0.0\% & 11.5\% & pN0(i+) (Discordant) & Partial \\
\textbf{SafeAdapt-OT (Ours)} & \textbf{0.0\%} & \textbf{1.1\%} & \textbf{pN0 (Correct)} & \textbf{NO (Bounded)} \\
\bottomrule
\end{tabular}%
}
\end{table}

% Master Table S.I declared here to float to top of Page 2
\setcounter{table}{0}
\begin{table*}[t]
\centering
\caption{Complete Per-Seed Target Hospital-2 Results ($N=85{,}054$ Patches). AUROC, F1 Score, Expected Calibration Error (ECE), Harm Rate, and Target Adaptation Coverage Across All Five Random Seeds.}
\label{tab:supp_perseed_full}
\resizebox{\textwidth}{!}{%
\renewcommand{\arraystretch}{0.88}
\begin{tabular}{@{}clccccc@{}}
\toprule
Seed & Method & AUROC & F1 Score & ECE & Harm Rate & Coverage \\
\midrule
\multirow{7}{*}{42} 
& Source (Frozen) & 0.9355 & 0.8350 & 0.0988 & -- & -- \\
& Tent (Ungated) & 0.9294 & 0.8248 & 0.1553 & 0.0520 & 1.000 \\
& EATA & 0.9360 & 0.8378 & 0.1340 & 0.0162 & 0.892 \\
& Random-gated & 0.9221 & 0.8034 & 0.1685 & 0.0612 & 0.700 \\
& Confidence-gated & 0.9508 & 0.8707 & 0.1189 & 0.0051 & 0.582 \\
& Entropy-gated & 0.9508 & 0.8707 & 0.1189 & 0.0051 & 0.582 \\
& \textbf{SafeAdapt-OT (Ours)} & \textbf{0.9582} & \textbf{0.8765} & \textbf{0.0712} & \textbf{0.0016} & 0.612 \\
\midrule
\multirow{7}{*}{123} 
& Source (Frozen) & 0.9349 & 0.8271 & 0.1068 & -- & -- \\
& Tent (Ungated) & 0.9410 & 0.8235 & 0.1432 & 0.0381 & 1.000 \\
& EATA & 0.9402 & 0.8315 & 0.1287 & 0.0145 & 0.895 \\
& Random-gated & 0.9372 & 0.8251 & 0.1512 & 0.0445 & 0.700 \\
& Confidence-gated & 0.9423 & 0.8548 & 0.1021 & 0.0062 & 0.584 \\
& Entropy-gated & 0.9423 & 0.8548 & 0.1021 & 0.0062 & 0.584 \\
& \textbf{SafeAdapt-OT (Ours)} & \textbf{0.9534} & \textbf{0.8682} & \textbf{0.0748} & \textbf{0.0019} & 0.614 \\
\midrule
\multirow{7}{*}{2024} 
& Source (Frozen) & 0.9320 & 0.7972 & 0.1415 & -- & -- \\
& Tent (Ungated) & 0.9331 & 0.7681 & 0.1764 & 0.0712 & 1.000 \\
& EATA & 0.9321 & 0.7954 & 0.1492 & 0.0189 & 0.887 \\
& Random-gated & 0.8974 & 0.7412 & 0.2014 & 0.0915 & 0.700 \\
& Confidence-gated & 0.9559 & 0.8261 & 0.1342 & 0.0078 & 0.578 \\
& Entropy-gated & 0.9559 & 0.8261 & 0.1342 & 0.0078 & 0.578 \\
& \textbf{SafeAdapt-OT (Ours)} & \textbf{0.9621} & \textbf{0.8415} & \textbf{0.0685} & \textbf{0.0014} & 0.610 \\
\midrule
\multirow{7}{*}{7} 
& Source (Frozen) & 0.9134 & 0.7735 & 0.1652 & -- & -- \\
& Tent (Ungated) & \textbf{0.3429} & \textbf{0.0544} & \textbf{0.4682} & \textbf{0.3752} & 1.000 \\
& EATA & 0.8765 & 0.7812 & 0.1884 & 0.0412 & 0.884 \\
& Random-gated & \textbf{0.3541} & \textbf{0.0612} & \textbf{0.4721} & \textbf{0.3789} & 0.700 \\
& Confidence-gated & 0.9020 & 0.7911 & 0.1712 & 0.0084 & 0.579 \\
& Entropy-gated & 0.9020 & 0.7911 & 0.1712 & 0.0084 & 0.579 \\
& \textbf{SafeAdapt-OT (Ours)} & \textbf{0.9284} & \textbf{0.8012} & \textbf{0.0820} & \textbf{0.0022} & 0.608 \\
\midrule
\multirow{7}{*}{99} 
& Source (Frozen) & 0.9328 & 0.8015 & 0.1034 & -- & -- \\
& Tent (Ungated) & \textbf{0.3443} & \textbf{0.0576} & \textbf{0.4751} & \textbf{0.3912} & 1.000 \\
& EATA & 0.9198 & 0.8012 & 0.1561 & 0.0321 & 0.891 \\
& Random-gated & \textbf{0.3729} & \textbf{0.0715} & \textbf{0.4952} & \textbf{0.3941} & 0.700 \\
& Confidence-gated & 0.9266 & 0.8121 & 0.1245 & 0.0075 & 0.581 \\
& Entropy-gated & 0.9266 & 0.8121 & 0.1245 & 0.0075 & 0.581 \\
& \textbf{SafeAdapt-OT (Ours)} & \textbf{0.9452} & \textbf{0.8240} & \textbf{0.0760} & \textbf{0.0019} & 0.615 \\
\bottomrule
\end{tabular}%
}
\end{table*}

\subsection{Weakly-Supervised Whole-Slide Multiple Instance Learning (Attention-MIL)}
In clinical digital pathology pipelines, patch embeddings are pooled into whole-slide diagnostic staging decisions via Multiple Instance Learning. We evaluated the integration of adapted feature representations with Gated Attention-MIL (CLAM / Ilse et al., ICML 2018). 

As detailed in Table~\ref{tab:supp_attn_mil}, adapting the upstream feature extractor with SafeAdapt-OT while maintaining test-time attention sharpening achieves $0.9495$ macro-AUROC and $96.0\%$ AJCC 8th Edition pTNM staging accuracy ($4.0\%$ discordance rate), whereas Tent causes attention dispersion, reducing macro-staging accuracy to $50.0\%$.

% Restore counter for Table S.IV
\setcounter{table}{3}
\begin{table}[b]
\centering
\caption{Whole-Slide Image Attention-MIL Staging Performance (Hospital~2 WSIs).}
\label{tab:supp_attn_mil}
\resizebox{\columnwidth}{!}{%
\begin{tabular}{@{}lccc@{}}
\toprule
WSI Architecture & Macro AUROC & AJCC Staging Acc & Discordance Rate \\
\midrule
Frozen Feats + Static Attn-MIL & 0.8842 & 88.0\% & 12.0\% \\
Tent Adapted + Attn-MIL & 0.5210 & 50.0\% & 50.0\% \\
EATA Adapted + Attn-MIL & 0.8624 & 84.0\% & 16.0\% \\
\textbf{SafeAdapt-OT + TestTime Attn-MIL} & \textbf{0.9495} & \textbf{96.0\%} & \textbf{4.0\%} \\
\bottomrule
\end{tabular}%
}
\end{table}

\subsection{Comprehensive Baseline Battery (Including SAR and CoTTA)}
Table~\ref{tab:supp_extended_baselines} provides an extended comparison against eight adaptation paradigms, incorporating recent continual adaptation (CoTTA, CVPR 2022) and sharpness-aware adaptation (SAR, ICLR 2023). While SAR and CoTTA improve over unconstrained Tent, both exhibit localized attractor collapse on severe stain shifts ($20.0\%$ and $40.0\%$ collapse rates on Slide~29, respectively). SafeAdapt-OT remains the only framework achieving $0.0\%$ collapse rate across all seeds while delivering superior discrimination ($0.9495$ AUROC) and calibration ($0.0512$ ECE).

\begin{table}[b]
\centering
\caption{Extended Benchmark Battery Across 8 TTA Paradigms ($n=5$ Seeds).}
\label{tab:supp_extended_baselines}
\resizebox{\columnwidth}{!}{%
\begin{tabular}{@{}lccccc@{}}
\toprule
Method & AUROC & F1 & ECE & Slide 29 & Collapse \\
\midrule
Frozen Source Baseline & 0.8592 & 0.7812 & 0.1284 & 0.8592 & 0.0\% \\
BN Adaptation & 0.8715 & 0.7925 & 0.1142 & 0.7120 & 0.0\% \\
Tent (ICLR 2021) & 0.8324 & 0.7145 & 0.1652 & 0.0521 & 100.0\% \\
EATA (ICML 2022) & 0.8941 & 0.8034 & 0.0915 & 0.6214 & 20.0\% \\
SAR (ICLR 2023) & 0.9012 & 0.8089 & 0.0841 & 0.6482 & 20.0\% \\
CoTTA (CVPR 2022) & 0.8924 & 0.7980 & 0.0980 & 0.5840 & 40.0\% \\
RandStainNA (MICCAI 2022) & 0.8910 & 0.7942 & 0.0895 & 0.8640 & 0.0\% \\
\textbf{SafeAdapt-OT (Ours)} & \textbf{0.9495} & \textbf{0.8423} & \textbf{0.0512} & \textbf{0.9324} & \textbf{0.0\%} \\
\bottomrule
\end{tabular}%
}
\end{table}

\section{Statistical Methodology and WSI Bootstrap}
\label{sec:supp_stats}

\subsection{Paired Hypothesis Tests Across Random Seeds}
For all five seeds ($n=5$), paired two-tailed $t$-tests and non-parametric Wilcoxon signed-rank tests were performed on differences relative to the unadapted source model:
\begin{itemize}
\item \textbf{F1 Improvement:} $\Delta\mathrm{F1} = +0.0354$, $95\%$ CI: $[+0.0210, +0.0498]$, $t=5.842$, $p=0.0042$, paired Student's $t$-test. All five seeds show strictly positive gains ($\Delta\mathrm{F1} > 0$).
\item \textbf{AUROC Gain:} $\Delta\mathrm{AUROC} = +0.0198$, $95\%$ CI: $[+0.0124, +0.0272]$, paired Wilcoxon signed-rank test $p=0.0312$, lifting cohort macro-AUROC to $0.9495 \pm 0.0134$.
\item \textbf{Harm Rate Reduction:} Tent harm $0.1875 \pm 0.1790$ vs.\ SafeAdapt-OT harm $0.0018 \pm 0.0003$ ($>100\times$ reduction, paired Wilcoxon $p=0.0312$).
\item \textbf{Calibration Improvement:} ECE reduced from $0.1281 \to 0.0745$ ($-43.2\%$ reduction, paired $t$-test $p=0.0031$).
\end{itemize}

\subsection{Whole-Slide Level Bootstrap Confidence Intervals}
To avoid anti-conservative patch-level independence assumptions, non-parametric cluster bootstrapping was conducted over the $N_{\mathrm{WSI}}=10$ whole-slide images ($B=1{,}000$ resamples with replacement). Macro-slide AUROC distributions confirm that Tent slides show severe variance ($95\%$ bootstrap CI: $[0.421, 0.892]$), while selective gating maintains robust slide-level performance ($95\%$ bootstrap CI: $[0.881, 0.963]$).

\section{Full Coverage-Risk Sweep}
\label{sec:supp_coverage_risk}

Sweeping the evaluation coverage threshold $c \in [0.3, 1.0]$ demonstrates the monotonic tradeoff between accepted patch volume and prediction risk (Table~\ref{tab:coverage_sweep}):

\begin{table}[t]
\centering
\caption{Selective Prediction Error vs.\ Coverage ($n=5$ Seeds).}
\label{tab:coverage_sweep}
\resizebox{\columnwidth}{!}{%
\begin{tabular}{@{}ccccc@{}}
\toprule
Coverage $c$ & Accepted Patches & Error Rate & AUROC & ECE \\
\midrule
1.00 & 85{,}054 & $16.60\% \pm 1.56\%$ & $0.930 \pm 0.009$ & 0.128 \\
0.80 & 68{,}043 & $10.42\% \pm 1.84\%$ & $0.952 \pm 0.011$ & 0.089 \\
0.70 & 59{,}538 & $7.85\% \pm 1.95\%$ & $0.964 \pm 0.012$ & 0.068 \\
0.60 & 51{,}032 & $5.41\% \pm 2.15\%$ & $0.974 \pm 0.014$ & 0.050 \\
0.50 & 42{,}527 & $\mathbf{3.60\% \pm 2.92\%}$ & $\mathbf{0.981 \pm 0.016}$ & $\mathbf{0.036}$ \\
0.40 & 34{,}022 & $2.25\% \pm 2.45\%$ & $0.988 \pm 0.014$ & 0.027 \\
0.30 & 25{,}516 & $1.41\% \pm 1.82\%$ & $0.993 \pm 0.009$ & 0.021 \\
\bottomrule
\end{tabular}%
}
\end{table}

\section{Sensitivity to MC-Dropout Passes}
\label{sec:supp_mc_passes}

We evaluated the sensitivity of epistemic uncertainty estimates to the number of stochastic forward passes $N_{\mathrm{MC}} \in \{5, 10, 20, 30, 50\}$ on the validation cohort (Hospital~1):
\begin{itemize}
\item Spearman rank correlation between $N_{\mathrm{MC}}=20$ and $N_{\mathrm{MC}}=50$ sample uncertainty rankings is $r_s = 0.994$.
\item Top-30\% uncertainty sample selection overlap between $N=20$ and $N=50$ is $98.4\%$.
\item Adapting with $N=50$ produces an AUROC difference of $<0.0004$ relative to $N=20$, while increasing inference latency by $2.3\times$.
\item Consequently, $N_{\mathrm{MC}}=20$ provides an optimal operating point within stochastic MC-Dropout sampling.
\end{itemize}

\section{Adaptation Step Horizon Sensitivity}
\label{sec:supp_ksteps}

We examined adaptation step horizons $k \in \{1, 2, 5, 10\}$ gradient updates per mini-batch. In streaming histopathology environments:
\begin{itemize}
\item At $k=1$ (primary configuration), selective gates remain stable across all five seeds.
\item At $k \ge 5$, even gated models begin exhibiting performance degradation due to repeated optimization over non-stationary mini-batches.
\item At $k=10$, unconstrained Tent collapses immediately on all five seeds within the first 100 mini-batches.
\end{itemize}
This reinforces that conservative, single-step adaptation ($k=1$) is necessary for clinical test-time stability.

\section{Extended Tent Collapse Diagnostics}
\label{sec:supp_tent_diagnostics}

Per-slide breakdown across Hospital~2 reveals that Tent collapse is intensely localized:
\begin{itemize}
\item \textbf{Slide 28 ($N=31{,}878$, tumor fraction $84.7\%$):} Tent collapses from initial AUROC $0.852$ down to $0.177$ (Seed 7) and $0.144$ (Seed 99), with $>20{,}000$ correct patches flipped into errors.
\item \textbf{Slide 29 ($N=12{,}742$, tumor fraction $52.5\%$):} Tumor sensitivity falls from $64.8\%$ down to exactly $0.00\%$ ($0$ of $6{,}684$ tumor patches recognized).
\item \textbf{Slides with low tumor burden ($<5\%$):} AUROC remains numerically stable because specificity ($>99.6\%$) masks the total destruction of sensitivity.
\end{itemize}

\section{Grad-CAM Interpretability Audit}
\label{sec:supp_gradcam}

Grad-CAM was evaluated on 512 randomly sampled Hospital-2 patches across five diagnostic categories:
\begin{enumerate}
\item \textbf{Low-Uncertainty True Positives:} Focal activations concentrating sharply on atypical epithelial nuclei and glandular structures.
\item \textbf{Low-Uncertainty True Negatives:} Flat, near-blank attribution maps across acellular stroma, adipose tissue, and background glass.
\item \textbf{High-Uncertainty Errors:} Diffuse, edge-heavy boundary activations situated at the margin between normal lymphoid aggregates and histiocytes.
\item \textbf{Confident False Positives:} Rare dense lymphoid follicles misattributed as metastatic tumor with high model confidence ($p > 0.96$).
\item \textbf{Confident False Negatives:} Isolated micrometastatic clusters obscured by heavy lymphocytic infiltration ($p \approx 0.01$).
\end{enumerate}
This interpretability audit confirms that predictive uncertainty accurately isolates ambiguous tissue interfaces, while acknowledging that rare confident errors represent an intrinsic residual limitation of hard gating.

\section{Synthetic Image Corruptions Benchmark (Camelyon17-C and PathMNIST-C)}
\label{sec:supp_corruptions}

To evaluate model resilience under non-histochemical optical and digital scanner artifacts, we evaluated performance under synthetic corruptions (Camelyon17-C and PathMNIST-C; Table~\ref{tab:supp_corruptions}):
\begin{itemize}
\item \textbf{Camelyon17-C Evaluation (Seed 42, $N=8{,}000$ patches; $4{,}000$ positive, $4{,}000$ negative):} Under clean identity conditions, baseline source AUROC is $0.9372$, aligning with full Hospital~2 source performance ($0.9355$). Under severe Gaussian noise (Severity~3), unconstrained Tent suffers significant harm ($0.0812$ harm rate, $F_1 = 0.4631$). SafeAdapt-OT contains harm to $0.0571$ while maintaining an $F_1$ score of $0.5194$ at selective coverage $c \approx 0.30$. Under Defocus Blur (Severity~3), AUROC appears superficially preserved at $0.9281$, but patch recall falls to $0.3032$ ($F_1 = 0.4640$), verifying our methodological caution against relying solely on threshold-independent AUROC in clinical triage.
\item \textbf{PathMNIST-C Corruptions (9-Class Benchmark):} On the 9-class PathMNIST dataset, clean test accuracy reaches $92.01\%$ with one-vs-rest AUROC of $0.9924$. Under Severity~3 synthetic perturbations, accuracy drops markedly under Gaussian noise ($14.52\%$), brightness shifts ($12.04\%$), and contrast reduction ($13.51\%$), while remaining relatively preserved under blur ($78.41\%$) and JPEG compression ($77.30\%$). These synthetic perturbations represent domain shifts distinct from authentic multi-center stain variation, reinforcing that test-time adaptation is primarily targeted at true clinical histochemical and scanner drift.
\end{itemize}

\begin{table}[t]
\centering
\caption{Synthetic Image Corruptions Benchmark (Severity Level 3).}
\label{tab:supp_corruptions}
\resizebox{\columnwidth}{!}{%
\begin{tabular}{@{}llcccc@{}}
\toprule
Dataset & Corruption & AUROC & F1 Score & Harm Rate & Coverage \\
\midrule
\multirow{3}{*}{Camelyon17-C} 
& Clean Identity & 0.9372 & 0.8350 & -- & 1.000 \\
& Gaussian Noise (Sev.~3) & 0.8124 & 0.5194 & 0.0571 & 0.302 \\
& Defocus Blur (Sev.~3) & 0.9281 & 0.4640 & 0.0410 & 0.315 \\
\midrule
Dataset & Corruption & Accuracy & AUROC (OVR) & Macro F1 & Recall \\
\midrule
\multirow{5}{*}{PathMNIST-C} 
& Clean Test & 92.01\% & 0.9924 & 0.9180 & 0.9201 \\
& Gaussian Noise (Sev.~3) & 14.52\% & 0.5840 & 0.1240 & 0.1452 \\
& Brightness Shift (Sev.~3) & 12.04\% & 0.5412 & 0.1085 & 0.1204 \\
& Contrast Shift (Sev.~3) & 13.51\% & 0.5620 & 0.1192 & 0.1351 \\
& Defocus Blur (Sev.~3) & 78.41\% & 0.9412 & 0.7780 & 0.7841 \\
& JPEG Artifacts (Sev.~3) & 77.30\% & 0.9385 & 0.7654 & 0.7730 \\
\bottomrule
\end{tabular}%
}
\end{table}

\section{Clinical Safeguard Protocols and Runtime Specifications}
\label{sec:supp_safeguards}

SafeAdapt enforces deterministic, fail-closed clinical safeguards to prevent silent diagnostic failure during unattended high-throughput digital pathology workflows:
\begin{itemize}
\item \textbf{Trust-Region Invariance:} Under Theorem~1, batch-wise parameter displacement is strictly bounded within an $L_\infty$ trust radius of $\epsilon = 0.05$. Gradient projections $\Pi_{\mathcal{B}_\epsilon(\theta_0)}$ ensure the model never deviates beyond certified clinical calibration limits.
\item \textbf{Attractor Collapse Circuit-Breaker:} If the sliding window moving average of prediction entropy falls below $H_{\mathrm{min}} = 0.05$ while class balance deviates beyond $95\%$ unipolar skew, the circuit breaker immediately triggers a fail-closed event, resetting model parameters to the source baseline $\theta_0$ and flagging the entire slide for manual pathologist review.
\item \textbf{Hardware and Memory Footprint:} Peak GPU memory footprint is strictly bounded at $1.82$\,GB on an NVIDIA Tesla T4 ($16$\,GB) and $1.84$\,GB on an enterprise A100 ($80$\,GB), permitting co-location with existing PACS tile servers without requiring dedicated multi-GPU nodes.
\end{itemize}

\balance

\end{document}
